\clearpage\thispagestyle{empty}
\pdfbookmark[0]{Quellen zur Übersicht über die Bahnbestimmung (1809)}{summarische1809-sources}
\begin{center}\textbf{Quellen zur Übersicht über die Bahnbestimmung (1809)}\end{center}
\noindent Diese KI-Übertragung in heutiges Deutsch folgt Carl Friedrich Gauß' deutschem Text in den \href{https://gdz.sub.uni-goettingen.de/id/PPN236007467}{\textit{Werken}, Band~VI} (Göttingen, 1874), S.~148--165. Sie enthält alle neun Artikel, die unterzeichnete einleitende Anmerkung von Bernhard von Lindenau und die vier Anmerkungen von Gauß. Der \href{https://www.digitale-sammlungen.de/en/view/bsb10538610?page=203}{Erstdruck} erschien als Beitrag XVII im Septemberheft 1809 der von Franz Xaver von Zach herausgegebenen \textit{Monatlichen Correspondenz zur Beförderung der Erd- und Himmels-Kunde}, Band~20, S.~197--224.
\par\smallskip\noindent Gauß verweist auf \href{https://www.e-rara.ch/zut/content/pageview/449839}{S.~45, §38} von Wilhelm Olbers' \textit{Abhandlung über die leichteste und bequemste Methode die Bahn eines Cometen aus einigen Beobachtungen zu berechnen} (Weimar, 1797). Zu den Integrationsformeln siehe den Nachtrag zu Roger Cotes' \textit{De methodo differentiali Newtoniana} in der von Robert Smith herausgegebenen und erweiterten \textit{Harmonia mensurarum} (Cambridge, 1722), \href{https://www.e-rara.ch/zut/content/pageview/1256011}{S.~32--33}. Der Erstdruck verweist auf \href{https://www.e-rara.ch/zut/content/pageview/142894}{S.~87} von Gauß' \textit{Theoria motus corporum coelestium in sectionibus conicis solem ambientium} (Hamburg, 1809). Dort beginnt Artikel~83, auf den die Werke verweisen.
\par\smallskip\noindent Eine eigene \href{https://www2.karlin.mff.cuni.cz/~tuma/Aplikace21/ALL_Gauss_1809_Summarische.pdf}{englische Fassung} ohne Angabe der übersetzenden Person ist über Jiří Tůmas Kurs Aplikace21 an der Karls-Universität zugänglich.
\par\smallskip\noindent Übertragung und Satz wurden mit KI erstellt.
\clearpage
